\subsection{Gram Type Classification}
% \noindent using the three envelopes ....
% - Results: coverage by class, set size distribution, singleton accuracy
% - analysis of hard cases: ambiguous gram-positive phages

% \subsubsection{Experimental Setup}
% \noindent - Data: split-0 phages only (80\% calibration, 20\% test), unknown host types excluded\\
% - Parameters: $\alpha = 0.1 , M=200$ directions (radial),
% M=20 bins (strip), 10 independent runs to estimate variance\\
% - The NC transformation used: keep scores as-is for one class, flip with 1 - s for the other 
% % (motivated by the empirical score semantics described in Chapter 4)

% \subsubsection{Results}
% \noindent - table: overall coverage, per-class coverage, average set size, empty set rate, size-2 rate, singleton accuracy std across 10 runs for all three methods\\
% - Bar charts: set size distribution, overall and per-class coverage


% \subsubsection{Analysis of Hard Cases}

% The minority class consistently has lower coverage than the majority class before the class-conditional fix — show this explicitly with a before/after comparison
% Identify the ambiguous phages: those whose mean raw score falls in the range [0.3,0.7][0.3, 0.7]
% [0.3,0.7] — these are phages for which the underlying classifiers are genuinely uncertain
% Show that coverage on typical phages of each class is much higher than coverage on the ambiguous subset — this is a property of the data, not a flaw in the conformal framework
% Discuss variance across runs: radial method has higher variance than strip and collapsed due to the random direction sampling